\documentclass[11pt]{article}

\usepackage[margin=1in]{geometry}
\usepackage{microtype}
\usepackage{amsmath}
\usepackage{graphicx}
\usepackage{booktabs}
\usepackage[hidelinks]{hyperref}

\graphicspath{{figures/}}

\newcommand{\natex}{\texttt{natex}}

\title{The Question Is the Treatment: A Known-Cutoff Regression Discontinuity
at the BTOS AI-Question Rewording}
\author{Hauke Hillebrandt\thanks{University College London. Email:
\texttt{ucjthhi@ucl.ac.uk}. This paper and the underlying \natex{} software were
prepared with substantial assistance from Anthropic's Claude models; the author
reviewed the analyses and text and is responsible for all remaining errors.}}
\date{July 2026 \\ \small analysis pass of record: 2026-07, \natex{} v0.2.0, seed 0}

\begin{document}

\maketitle

\begin{abstract}
\noindent Beginning with wave 202524 (reference period from 3 November 2025;
announced 3 December 2025), the US Census Bureau's Business Trends and Outlook
Survey (BTOS) reworded its headline AI question from AI use ``in producing
goods or services'' to AI use ``in any of its business functions.'' Treating
the refresh as a known measurement cutoff in a regression-discontinuity-in-time
design on the spliced 71-wave national series, the rewording moved measured US
business AI use by $+6.04$ percentage points (HC1 $z=22.9$; HAC $z=16.3$) ---
a ratio of 1.553 over the old-wording level at the cutoff. The jump is $4.7\times$
the largest of 49 placebo-cutoff estimates (randomization $p=0.020$, the $1/50$
floor), stable at $6.0$--$7.5$ pp across bandwidth, weighting, curvature, and
HAC variants, and a blind LoRD3 scan re-localizes it: the top two data-driven
candidate cutoffs exactly flank the known one (honest 2SLS $+6.74$, 95\% CI
$[6.01, 7.48]$). True adoption growth cannot produce it: the fitted pre-trend
predicts $+0.66$ pp across the 50-day shutdown gap, and doubling that slope
still leaves $\geq 5.38$ pp. Nulls are reported as nulls: the blind-scan
randomization test is underpowered at $n=71$ ($p=0.14$), the McCrary density
test is inapplicable on a uniform time grid (NaN), a December--January
seasonal term is insignificant ($0.29 \pm 0.24$), and a post-cutoff slope
steepening ($+2.17$ pp/yr, conventional se $0.56$) is secondary and suggestive
only. The estimate calibrates a splice factor --- additive $+6.04$ pp or ratio
$1.553$ --- for carrying the pre-refresh series past December 2025, with the
caveat that the ``wording effect'' is bundled with a simultaneous sample
rotation and any shutdown-adjacent change in response composition.
\end{abstract}

\section{Introduction}\label{sec:intro}

Since September 2023 the Census Bureau's Business Trends and Outlook Survey
(BTOS) has asked a probability sample of US firms whether they used artificial
intelligence in the last two weeks; \cite{bonney2024} introduce these data,
which have become the flagship high-frequency gauge of US business AI adoption
--- the headline series in the Federal Reserve's monitoring of the AI economy
\cite{allen2026}. On 3 December 2025 the Census Bureau announced that,
beginning with survey cycle 2 of sample year 2026, the AI questions were
reworded: the core use question now asks about AI ``in any of its business
functions'' rather than ``in producing goods or services,'' explicitly to
reduce respondent under-reporting \cite{census2025}. Measured adoption moved
immediately: the last old-wording wave reads 10.0\%, the first new-wording wave
17.3\%. Commentary quickly flagged the break --- \cite{bick2026} show that how
the question is asked drives large gaps between firm AI-adoption estimates,
and \cite{allen2026} notes the redesign when reporting the $\sim$18\%
year-end-2025 level --- and Census researchers now exploit the expanded
business-function detail the redesign introduced \cite{bonney2026}. What the
commentary does not provide is a formal estimate of the rewording effect with
falsification-calibrated uncertainty. That estimate is the number any user of
the series needs in order to splice the pre- and post-refresh segments into
one panel.

Instrument refreshes breaking measured series is a classic survey-methods
problem. The canonical precedent is the 1994 Current Population Survey
redesign, whose effects were quantified with a dedicated parallel survey run
under both instruments and converted into adjustment factors for splicing the
historical series \cite{polivka1998}. The BTOS refresh has no parallel
overlap sample: the old question stopped, the new one started, and a federal
government shutdown sits between them. The splice factor must therefore be
estimated from the time series alone --- a known-cutoff
regression-discontinuity-in-time (RDiT) problem \cite{hausman2018}, here in
its cleanest form: the cutoff is announced, mechanical, and datable to a
single wave, and the ``treatment'' is the measurement instrument itself, so
the usual manipulation and anticipation concerns are moot while the RDiT
time-series concerns (serial correlation, trend extrapolation) remain and are
addressed directly. We estimate the jump by local linear regression on each
side of the cutoff with robust and HAC inference \cite{calonico2014,
hausman2018}, calibrate it against a 49-point placebo-cutoff grid, bound the
adoption-growth confound, and additionally ask the \natex{} toolkit
\cite{natex} to find the discontinuity \emph{blind}, via the LoRD3 local-RD
scan of \cite{herlands2018}, plus a secondary regression-kink test in the
conventions of \cite{card2015} and \cite{bockerman2025}.

\section{Data}\label{sec:data}

The outcome is the national share of firms answering ``Yes'' to BTOS Question
7, in percentage points, with published survey standard errors. Two vintages
are required because current Census downloads carry only the reworded
question, with all earlier waves blanked. The old-wording segment (``\ldots
did this business use Artificial Intelligence (AI) in producing goods or
services?'') comes from a frozen pre-refresh AI-core workbook: 54 biweekly
waves, 202319--202520, rising 3.7\% to 10.0\%. The new-wording segment
(``\ldots in any of its business functions?'') comes from a fresh
\texttt{National.xlsx} download (census.gov, 2026-07-18, 82~kB, verified):
17 waves, 202524--202614, rising 17.3\% to 21.7\%. The spliced panel has
$n=71$ waves indexed by $t$, the decimal-year reference-period midpoint.

Three timeline facts matter. First, waves 202521--202523 were never collected
--- the October--November 2025 federal government shutdown --- so the design
has a built-in donut that is a no-data gap, not censoring: 50 days separate
the last pre-cutoff reference midpoint (2025-09-14) from the cutoff. Second,
the cutoff is $t_c = 2025.8384$ (2025-11-03, the reference-period start of
wave 202524, the first wave fielded with the reworded question; published
2025-12-04). No reference midpoints fall in (2025-09-14, 2025-11-09), so any
cutoff in that interval yields the identical estimation sample. Third, wave
202524 also begins Sample Year 4 (a scheduled full sample rotation), so the
refresh bundles the wording change with a panel rotation --- a caveat carried
throughout.

\section{Design and Methods}\label{sec:methods}

\paragraph{Estimator.} The headline specification fits SE-weighted
($w=1/\mathrm{se}^2$) local linear regressions of the AI-use share on
$t - t_c$ separately on each side of the cutoff within bandwidth $h=1.0$
years ($n_{\text{pre}}=23$, $n_{\text{post}}=17$); the jump $\tau$ is the
post-side intercept minus the pre-side extrapolation at $t_c$. Inference is
HC1 per side combined by independence, plus a pooled HAC (Bartlett kernel, 4
and 8 lags) variant for serial correlation. Sensitivity covers bandwidths
$h \in \{0.5, 0.75, 1.0, 2.3\}$, unweighted fits, and pre-side quadratics.

\paragraph{Identification and falsification.} In a measurement RDD the
treatment is the instrument refresh, so the only smooth confound is true
adoption growth loading into the 50-day extrapolation window; we bound it by
the fitted pre-slope and report the estimate under a doubled gap slope.
Calibration is design-based: the identical estimator is run at all 49
feasible interior placebo cutoffs, giving a randomization $p$-value with
floor $1/50$. A blind check uses \texttt{natex discover} (LoRD3, seed 0,
$k=12$, $q=99$): the scan searches for local discontinuities without being
told the cutoff, then honest (split-sample) 2SLS re-estimates the effect at
the discovered boundary \cite{herlands2018,natex}. The McCrary density test
\cite{mccrary2008} returns NaN on a uniform time grid --- inapplicable in
RDiT, reported as an outcome, not patched over. A secondary sharp
regression-kink test at $t_c$ (\texttt{natex kink}, right-minus-left slope
contrast, policy kink 1, $h=1.0$, conventional local-polynomial inference
without HAC) asks whether the trend also steepened \cite{card2015,
bockerman2025}.

\paragraph{Honest-inference notes, stated as run.} The pipeline's refusals
and artifacts are reported: the automated \texttt{natex survey} run refused
its rdd family (100-row floor, $n=71$), iv ($\geq 80$), did/sc (needs a
panel), and dee ($\geq 200$); its kink family auto-assigned the degenerate
outcome \texttt{run} (a linear transform of $t$), so that null
($p=0.79$, $\tau \approx 10^{-13}$) is meaningless and was superseded by the
declared-outcome kink run above. The LoRD3 validation battery passed
(covariate placebo Holm $p=0.24$); its blind-scan randomization test is
underpowered at $n=71$ ($p=0.14$) and is reported as the null it is.
Pre-fit residuals have AR(1) $=0.836$ and empirical wave noise $\sim$4$\times$
the published survey SEs (residual sd $0.614$ vs.\ mean reported se $0.151$);
the HC1/HAC empirical standard errors absorb this overdispersion, but HAC
with $n_{\text{post}}=17$ is approximate --- the jump is $\sim$16 HAC standard
errors, so no plausible correction overturns it. A December--January seasonal
coefficient is insignificant ($0.293 \pm 0.236$).

\section{Results}\label{sec:results}

\begin{figure}[t]
\centering
\includegraphics[width=\textwidth]{fig1}
\caption{(a) The spliced BTOS national AI-use series with SE-weighted local
linear fits within $h=1.0$ yr of the cutoff $t_c=2025.8384$ (2025-11-03,
vertical line); the dashed segment is the 50-day extrapolation across the
shutdown gap (shaded); the arrow marks the $+6.04$ pp jump. (b) The same
estimator at all 49 feasible placebo cutoffs (open circles) versus the true
cutoff (filled, with 95\% HAC confidence interval): the largest placebo is
$1.30$ pp, $4.7\times$ smaller than the true jump (randomization $p=0.020$).}
\label{fig:fig1}
\end{figure}

\paragraph{Headline jump and robustness.} At the known cutoff the rewording
jump is $\hat\tau = +6.039$ pp (HC1 se $0.263$, $z=22.9$; HAC(4) se $0.371$,
$z=16.3$; HAC(8) se $0.376$), against an old-wording level of $10.92$ pp at
the cutoff --- a ratio of $1.553$. Every variant in
Table~\ref{tab:rd} stays within $6.0$--$7.5$ pp with $z>12$: bandwidths
$0.5$/$0.75$/$2.3$ give $7.148$/$6.383$/$7.363$, the unweighted fit gives
$5.956$, and pre-side quadratics give $5.787$--$7.464$.

\paragraph{Placebo-cutoff calibration.} Across the 49 feasible interior
placebo cutoffs the same estimator yields a maximum $|\tau|$ of $1.298$ pp
(mean $0.515$); none reaches the true jump in $|\tau|$ or $|z|$, for a
randomization $p = 0.020$ --- the minimum achievable with 49 placebos
(Figure~\ref{fig:fig1}b). The true jump is $4.7\times$ the largest placebo.

\paragraph{Blind re-localization.} Without being told the cutoff, the LoRD3
scan's two highest-scoring candidate boundaries sit at $t=2025.7027$ and
$t=2025.8562$ --- exactly the last old-wording and first new-wording waves,
flanking $t_c$. Honest 2SLS at the discovered boundary gives $\tau = 6.743$
(se $0.375$, 95\% CI $[6.008, 7.478]$, no weak-instrument flag); the naive
full-sample Wald estimate is $8.06$ (se $0.206$). The scan-level
randomization test does not reach significance ($p=0.14$), as expected at
$n=71$; localization, not the scan $p$-value, is the evidential contribution.

\paragraph{Bounding the adoption confound.} The fitted pre-slope is
$4.89 \pm 0.31$ pp/yr, which predicts only $+0.66$ pp of true-adoption growth
across the 50-day gap. Even if the true gap slope were double the fitted
slope, the implied jump remains $\geq 5.376$ pp. A wording effect of zero
would require adoption to grow $\sim$9$\times$ faster during the shutdown
than in the preceding year --- and then revert, since the post-cutoff slope
($6.43$ pp/yr) is only modestly steeper.

\paragraph{Secondary kink.} The sharp RKD at $t_c$ estimates a slope change
of $+2.170$ pp/yr (conventional se $0.558$, no HAC): the new-wording series
trends steeper than the local old-wording trend. Given conventional
inference and the bundling caveats, this is suggestive only; it cautions
against assuming the wording effect is a pure level shift far from the
cutoff.

\paragraph{Splice calibration.} The fitted levels at the cutoff are $10.92$
pp (old wording) and $16.96$ pp (new wording): an additive splice of
$+6.04$ pp or a ratio splice of $1.553$. Over the observed post period the
two disagree increasingly far from the cutoff (a ratio splice scales the
steeper post trend); for level-tracking the additive factor is the direct RD
estimand, and the ratio is the natural choice if the wording effect scales
with the underlying level.

\begin{table}[t]
\centering
\small
\caption{RD estimates of the BTOS rewording jump at $t_c=2025.8384$.}
\label{tab:rd}
\begin{tabular}{lrrr}
\toprule
Specification & $\hat\tau$ & se & $z$ \\
\midrule
\multicolumn{4}{l}{\emph{Panel A: known-cutoff RD jump (pp), SE-weighted local linear}}\\
$h=1.0$, HC1 (headline)          & $+6.039$ & $0.263$ & $22.9$ \\
$h=1.0$, HAC Bartlett, 4 lags    & $+6.039$ & $0.371$ & $16.3$ \\
$h=1.0$, HAC Bartlett, 8 lags    & $+6.039$ & $0.376$ & $16.1$ \\
$h=0.5$                          & $+7.148$ & $0.352$ & $20.3$ \\
$h=0.75$                         & $+6.383$ & $0.275$ & $23.2$ \\
$h=2.3$ (full sample)            & $+7.363$ & $0.241$ & $30.6$ \\
$h=1.0$, unweighted              & $+5.956$ & $0.287$ & $20.8$ \\
$h=0.75$, pre-quadratic          & $+7.464$ & $0.481$ & $15.5$ \\
$h=1.0$, pre-quadratic           & $+7.051$ & $0.374$ & $18.9$ \\
$h=2.3$, pre-quadratic           & $+5.787$ & $0.282$ & $20.5$ \\
\midrule
\multicolumn{4}{l}{\emph{Panel B: blind discovery, placebos, and secondary tests}}\\
LoRD3 honest 2SLS (blind scan)   & $+6.743$ & $0.375$ & $18.0$ \\
LoRD3 naive Wald                 & $+8.063$ & $0.206$ & $39.2$ \\
Placebo cutoffs (49): max $|\tau|$ & $1.298$ & --- & --- \\
Placebo cutoffs (49): mean $|\tau|$ & $0.515$ & --- & --- \\
Sharp RKD slope change (pp/yr)   & $+2.170$ & $0.558$ & $3.9$ \\
\bottomrule
\end{tabular}

\medskip
\parbox{0.92\textwidth}{\footnotesize Notes: outcome is the national share
answering Yes to BTOS Question 7, percentage points; $n=23$ pre / $17$ post
waves at $h=1.0$ yr. HC1 per side combined by independence; HAC is a pooled
Bartlett-kernel variant. Placebo randomization $p=0.020$ ($0$ of $49$
placebos $\geq$ the true jump in $|\tau|$ or $|z|$; floor $1/50$). LoRD3:
\texttt{natex discover}, seed 0, $k=12$, $q=99$; scan randomization $p=0.14$
(underpowered at $n=71$); McCrary density test NaN (uniform time grid,
inapplicable). RKD: conventional local-polynomial inference, no HAC;
secondary.}
\end{table}

\section{Caveats and Conclusion}\label{sec:conclusion}

Four caveats bound the claim. First, the donut is a federal-shutdown no-data
gap (waves 202521--202523 were never collected), so any true adoption
acceleration inside the 50-day extrapolation window loads into $\hat\tau$;
the pre-slope bound caps this at $0.66$ pp under the fitted trend and leaves
$\tau \geq 5.38$ pp even at double that slope. Second, the questionnaire
refresh is bundled with a full sample rotation (Sample Year 4 begins at wave
202524) and follows the shutdown, so ``wording effect'' strictly means
wording $+$ panel-refresh $+$ shutdown-adjacent response-composition effect;
for the splice-factor purpose all three are measurement artifacts, but pure
wording attribution is not separable in this design. Third, RDiT caveats
apply \cite{hausman2018}: pre-fit residual AR(1) is $0.836$ and empirical
wave noise is $\sim$4$\times$ the published survey SEs; the HC1/HAC empirical
standard errors absorb this, but HAC with 17 post waves is approximate ---
the jump is $\sim$16 HAC standard errors, so no plausible correction
overturns it. Fourth, the automated-pipeline artifacts were triaged, not
hidden: the survey rdd family refused on its 100-row floor, the auto-assigned
degenerate kink outcome was superseded by the declared-outcome run, the
McCrary test is inapplicable by construction, and the blind-scan
randomization $p=0.14$ is an underpowered null.

The verdict is credible: a textbook known-cutoff measurement RDD. The
$+6.04$ pp jump at the 2025-11-03 refresh is $4.7\times$ the largest of 49
placebo jumps, survives every bandwidth, weighting, curvature, and HAC
variant within a $6.0$--$7.5$ pp band, is blindly re-localized by LoRD3 at
exactly the flanking waves, and cannot be true adoption. It identifies the
instrument-refresh effect and calibrates the splice --- additive $+6.04$ pp,
ratio $1.553$ --- required to extend any analysis built on the old-wording
BTOS AI series past December 2025. Substantively, the estimate implies that
roughly a third of the post-refresh measured adoption level is attributable
to asking a better question rather than to more firms using AI --- a caution
for any trend analysis that splices BTOS vintages naively, and a
quantification of the under-reporting the Census Bureau redesigned the
question to remove \cite{census2025,bick2026}.

\paragraph{Reproducibility.} All estimates were produced with \natex{}
v0.2.0 \cite{natex} at seed 0 from the frozen BTOS workbooks (published at
\url{https://www.census.gov/hfp/btos}; not committed to the repository).
Figure~\ref{fig:fig1} regenerates deterministically from the committed
\texttt{figures/make\_fig.py}, which asserts the headline estimate, the
refitted cutoff levels, and the placebo maximum against the numbers of
record before drawing.

\footnotesize
\begin{thebibliography}{9}
\setlength{\itemsep}{2pt}\setlength{\parskip}{0pt}

\bibitem{allen2026}
Allen, J.~S. (2026).
Monitoring AI adoption in the U.S. economy.
\emph{FEDS Notes}, Board of Governors of the Federal Reserve System, 3 April
2026.

\bibitem{bick2026}
Bick, A., Blandin, A., Deming, D., Fuchs-Sch\"undeln, N., and Jessen, J.
(2026).
Measuring AI adoption among firms: How you ask matters.
Federal Reserve Bank of St.\ Louis, \emph{On the Economy}, June 2026.

\bibitem{bockerman2025}
B\"ockerman, P., Jysm\"a, S., and Kanninen, O. (2025).
\emph{Difference-in-Kinks Design}.
IZA Discussion Paper No.~18313.
\url{https://docs.iza.org/dp18313.pdf}

\bibitem{bonney2024}
Bonney, K., Breaux, C., Buffington, C., Dinlersoz, E., Foster, L., Goldschlag,
N., Haltiwanger, J., Kroff, Z., and Savage, K. (2024).
Tracking firm use of AI in real time: A snapshot from the Business Trends and
Outlook Survey.
NBER Working Paper No.~32319.

\bibitem{bonney2026}
Bonney, K., Breaux, C., Dinlersoz, E., Foster, L., Haltiwanger, J., and
Pande, A. (2026).
The microstructure of AI diffusion: Evidence from firms, business functions,
and worker tasks.
Census Bureau CES Working Paper CES-WP-26-25; NBER Working Paper No.~35141.

\bibitem{calonico2014}
Calonico, S., Cattaneo, M.~D., and Titiunik, R. (2014).
Robust nonparametric confidence intervals for regression-discontinuity
designs.
\emph{Econometrica}, 82(6), 2295--2326.

\bibitem{card2015}
Card, D., Lee, D.~S., Pei, Z., and Weber, A. (2015).
Inference on causal effects in a generalized regression kink design.
\emph{Econometrica}, 83(6), 2453--2483.

\bibitem{census2025}
U.S. Census Bureau (2025).
\emph{Business Trends and Outlook Survey: AI core question updates}.
3 December 2025.
\url{https://www.census.gov/hfp/btos/downloads/AI\%20Question\%20Wording\%20Updates.pdf}

\bibitem{hausman2018}
Hausman, C., and Rapson, D.~S. (2018).
Regression discontinuity in time: Considerations for empirical applications.
\emph{Annual Review of Resource Economics}, 10, 533--552.

\bibitem{herlands2018}
Herlands, W., McFowland~III, E., Wilson, A.~G., and Neill, D.~B. (2018).
Automated local regression discontinuity design discovery.
In \emph{Proceedings of the 24th ACM SIGKDD International Conference on
Knowledge Discovery and Data Mining (KDD~'18)}.

\bibitem{mccrary2008}
McCrary, J. (2008).
Manipulation of the running variable in the regression discontinuity design:
A density test.
\emph{Journal of Econometrics}, 142(2), 698--714.

\bibitem{natex}
Hillebrandt, H. (2026).
\emph{natex: automated natural-experiment discovery and estimation} (version
0.2.0). Software.
\url{https://github.com/HaukeHillebrandt/natex}

\bibitem{polivka1998}
Polivka, A.~E., and Miller, S.~M. (1998).
The CPS after the redesign: Refocusing the economic lens.
In \emph{Labor Statistics Measurement Issues}, 249--289.
University of Chicago Press.

\end{thebibliography}

\end{document}
